(1) 容器内速度分布是 $f(\mathbf v)=(m/2\pi k_BT)^{3/2}e^{-m|\mathbf v|^2/(2k_BT)}$。通过孔的事件数带有向外法向速率权重 $nv_zf(\mathbf v)d^3v$（$v_z>0$），因为快速分子在相同时间扫过更大的孔前体积。

(2) 对半空间速度积分，单位面积粒子通量为 $\boxed{j_N=n\int_{v_z>0}v_zf d^3v=n\sqrt{k_BT/(2\pi m)}}$。事件条件下平均法向速度是 $\langle v_z\rangle_{\rm eff}=\langle v_z^2\rangle_{v_z>0}/\langle v_z\rangle_{v_z>0}=\sqrt{\pi k_BT/(2m)}$，高于普通半空间平均法向速率 $\sqrt{2k_BT/(\pi m)}$。

(3) 逸出通量加权的能量是 $\langle mv^2/2\rangle_{\rm eff}=2k_BT$。可用横向两个自由度各平均 $k_BT/2$，法向因速度加权而平均 $k_BT$ 直接相加。故单位面积能量流为 $\boxed{j_E=2k_BT,j_N}$，比体内每粒平均动能 $3k_BT/2$ 高。

(4) 若容器体积 $V$ 恒定，$N=nV$，则 $\dot N=-Aj_N=-(AN/V)\sqrt{k_BT/(2\pi m)}$。内部总能 $U=3Nk_BT/2$，能量平衡 $dU/dt=2k_BT\dot N$（逸出使 $\dot N<0$），因此 $\boxed{(3/2)Nk_B\dot T=(1/2)k_BT\dot N}$，即 $\dot T/T=(1/3)\dot N/N<0$。要求内部快速碰撞重新热化且孔足够小，不计内部势能和反向回流。
